\chapter{Differences from libroadrunner}
\label{ch:differences}

roadrunner\_nitro follows libroadrunner's behaviour, including its defaults, except where libroadrunner's
behaviour is wrong. Over the 1075 curated BioModels, it simulates every model libroadrunner simulates, and 45
that libroadrunner cannot. Time courses agree to a relative $10^{-5}$ on 959 models; where they differ, COPASI as
an independent referee puts both engines within $10^{-3}$ of its answer in most cases, with roadrunner\_nitro
the closer more often. The deliberate differences are:

\begin{center}
\begin{tabularx}{\textwidth}{L{0.36\textwidth}YY}
\toprule
Function & libroadrunner & roadrunner\_nitro \\
\midrule
\code{getRatesOfChange} & values in a different order from \code{getRatesOfChangeIds} & values match the ids \\
\code{setTimeCourseSelectionList} with \code{"time, cc(S1, k1)"} & splits inside \code{cc(...)} at the comma & splits only
  between selections \\
\code{loadSimulationSettings} & settings have no effect & start, duration, steps and tolerances are applied \\
\code{getFloatingSpecies\-Initial\-Concentration\-ByIndex} & a bad index returns success without a value & a bad
  index is an error \\
\code{getConfigurationXML}, \code{writeRRData} & declared but missing & implemented \\
\code{steadyState} & gives up when Newton's method fails from the current state & also finds steady states
  Newton's method alone misses \\
euler, rk4, rk45 integrators & ignore events after the start time & apply every event \\
\bottomrule
\end{tabularx}
\end{center}

In Python, the package is imported as \code{roadrunner\_nitro} rather than \code{roadrunner}, and it adds
\code{simulateBatch} (many time courses in parallel) and the \code{fit} module, which libroadrunner does not
have.
